
\subsection{Additional Details for Approximate K-Means}\label{app:kmeans}

\subsubsection{Proof for Error-bound Guarantees}


We restate the notations and propositions here for clarity. Let $X= \{ x_1, \ldots, x_n\}$ and $\tX = \{ \tilde{x_1}, \ldots, \tilde{x_n}\}$ denote a set of $n$ points where $x_i, \tx_i \in [0,1]$ and for all $i$, $|x_i - \txi| \leq \alpha x_i$ for some $\alpha \in (0,1)$. Let $M = \{\mu_1, \ldots, \mu_k\}$ be the cluster centers of $X$ and the loss of $M$ over inputs $X$ is defined as $\loss(M) = \frac{1}{n}\sum_{x \in X} \min_{\mu \in M} |x - \mu|$. The cluster centers $\tM$ and loss $\qloss$ over $\tX$ are defined similarly. The optimal loss $\loss(\oM)$ is the loss of the optimal cluster centers $M^{*} = \argmin_{M} \loss(M)$. The optimal loss $\qloss(\otM)$ and cluster centers $\otM$ over $\tX$ are defined similarly. Let $\mu_{M}(i) = \argmin_{\mu} |\mu - x_i|$ denote the cluster center in $M$ that is closest to the ith input $x_i$ ($\tmu_{\tM}(i)$ is defined similarly for $\txi$).



Note that the loss $\loss(\tM)$ over the original inputs $X$ with any clustering $\tM$ over $\tX$ is at most $\alpha$ greater than the loss $\qloss(\tM)$,

\[   
    \qloss(\tM) = \frac{1}{n} \sum_{i=1}^n |\tmu_{\tM}(i) - \txi|  \leq \frac{1}{n} \sum_{i=1}^n (|\tmu_{\tM}(i) - x_i| + \alpha x_i) 
\]

\[
 = \loss(\tM) + \frac{\alpha}{n} \sum_{i=1}^n x_i \leq \loss(\tM) + \alpha
\]

where the first inequality follows from triangle inequality and the second follows from the fact that $\frac{1}{n}\sum_{i=1}^n x_i \leq 1$. Using similar arguments, one can verify that $\loss(\tM) \leq \qloss(\tM) + \alpha$.

\begin{proposition}
Given any two sets $X, \tX$ of $n$ points as defined above such that $|x_i - \txi| \leq \alpha x_i$ $\forall i$. For any $k \geq 1$, the difference between the optimal loss of clusterings over $X$ and $\tX$ are bounded as, 

\begin{equation}
    |\loss(\oM) - \qloss(\otM)| \leq \alpha
\end{equation}
\end{proposition}

\begin{proof} We first show that $\qloss(\otM) \leq \loss(\oM) + \alpha$.
   \[
   \qloss(\otM) = \frac{1}{n} \sum_{i=1}^n |\tmui - \txi| \leq \frac{1}{n} \sum_{i=1}^n |\mui - \txi| 
   \]
   \begin{equation}\label{Eq:boundsize1}
   \leq \frac{1}{n} \left ( \sum_{i=1}^n |\mui - x_i| + \alpha \sum_{i=1}^n x_i \right) \leq \loss(\oM) + \alpha
   \end{equation}

In \cref{Eq:boundsize1}, the first inequality follows from the fact that $\otM$ is the optimal clustering over $\tX$ and hence replacing the cluster centers $\tmui$ with $\mui$ leads to a quantity which is greater than equal to $\qloss(\otM)$. The next inequality follows from triangle inequality and uses the fact that $|x_i - \txi| \leq \alpha x_i$. Using similar arguments, one can verify that $\loss(\oM) \leq \qloss(\otM)  + \alpha$. Combining the two equations, we have that $|\loss(\oM) - \qloss(\otM)| \leq \alpha$ which is the main statement of the proposition.

The proof for $\loss(\oM) \leq \qloss(\otM)  + \alpha$ follows similar arguments leading to the main statement of the result.


\end{proof}

\begin{remark} The results show that the difference between optimal loss of k-means clustering over inputs $X$ and its grouped counterpart $\tX$ is bounded by $\alpha$. Building upon this, we will show that when the clustering is performed with sample-weighted k-means++, the expected loss over $\tX$ is bounded by the optimal loss over $X$ and $\alpha^2$. 
\end{remark}


To keep notations close to \cite{kmeanspp}, we denote the loss over squared norm as $\phi(M) = \frac{1}{n}\sum_{x \in X} \min_{\mu \in M} |x - \mu|^2$. The loss $\phiq(M)$ is defined similarly over $\tX$: $\phiq(\tM) = \frac{1}{n}\sum_{\tx \in \tX} \min_{\tmu \in \tM} |\tx - \tmu|^2$. 

\begin{proposition}\label{prop2}
Given any two sets $X, \tX$ of $n$ points as defined above such that $|x_i - \txi| \leq \alpha x_i$ $\forall i$. For any $k \geq 1$, let $\wM$ be the clustering obtained by applying weighted kmeans++ (Algorithms \ref{algo:initialize} and \ref{algo:kmeans}) over the set $\tX$. Then,

\begin{equation}
    E[\phiq(\wM)] \leq 16 (\ln k + 2) (\phi(\oM) + \alpha^2)
\end{equation}

\end{proposition}
\begin{proof}
    We first upper bound the optimal loss $\phiq(\otM)$ with $\phi(\oM)$ and $\alpha$. See that

 \[
   \phiq(\otM) =  \frac{1}{n} \sum_{i=1}^n |\tmui - \txi|^2  \leq \frac{1}{n} \sum_{i=1}^n |\mui - \txi|^2
   \]
   \begin{equation}\label{Eq:squaredbound}
   \leq 2 \left ( \frac{1}{n}\sum_{i=1}^n |\mui - x_i|^2 + \frac{\alpha^2}{n} \sum_{i=1}^n x_i^2 \right)  \leq 2(\phi(\oM) + \alpha^2)
   \end{equation}

   Applying weighted-kmeans++ (Algorithms \ref{algo:initialize} and \ref{algo:kmeans}) over the set $\tX$ produces $\wM$. Using Theorem 1.1 in \cite{kmeanspp}, we have that $E[\phiq(\wM)] \leq 8 (\ln k + 2) (\phiq(\otM))$ and using Eq. \ref{Eq:squaredbound} to replace $\phiq(\otM)$, we obtain the main statement of the proposition.

\end{proof}

\begin{remark} The result implies that \cref{algo:initialize,algo:kmeans} is $O(\ln k)$-competitive with $\phi(\oM) + \alpha^2$ while being orders of magnitude faster in comparison with the vanilla k-means++ which is $O(\ln k)$-competitive with $\phi(\oM)$.
\end{remark}



We now analyze certain conditions where the clustering memberships of every point in the two cases ($X$ and $\tX$) would be identical. Let $C = \{X_1, \ldots, X_k\}$ be a clustering (or partition) of $X \subseteq [0, 1]$ with centers $M = \{\mu_1, \ldots, \mu_k\}$. For all $x \in X_i$, and $i \neq j$, $|x - \mu_i| \leq |x - \mu_j|$. Let the minimum distance between any two points in two different clusters in $C$ be the split of the clustering denoted as $split_C(X)$. Let the maximum distance between any two points in the same cluster be the width of the clustering denoted as $width_C(X)$.

\begin{definition}
   \cite{ben2014clustering} A clustering $C = \{X_1, \ldots, X_k\}$ of $X \subseteq [0, 1]$ is $\sigsep$ for $\sigma \geq 1$ if  $split_C(X) > \sigma \cdot width_C(X)$.
\end{definition}

In other words, any clustering $C$ is $\sigsep$ if the minimum distance between any two points in separate clusters is greater than the maximum distance between two points in the same cluster. 

We now show that if the optimal clustering over the set $X$ is $\sigsep$, then for small enough $\alpha$ depending on the value of $\sigma$, the optimal clustering over $X$ and $\tX$ will be the same. 

First note that, if a $\sigsep$ clustering $C$ exists for a given set of inputs, then only one such paritioning of the inputs. 

\begin{lemma}
    \cite{ackerman2009clusterability} If there exists a k-clustering $C$ of $X$, for $k \geq 2$, such that $C$ is $\sigsep$, then there is only one such partitioning of $X$.
\end{lemma}

The proof of the above Lemma can be found in \cite{ackerman2009clusterability}. We now show that if $\alpha < \frac{split_C(X) - width_C(X)}{4}$, then the same paritioning over $\tX$ is also $\sigsep$ and hence the optimal clustering for both $X$ and $\tX$ are identical.

Let $s_1$ and $s_2$ be the two points in $X$ such that the distance between them is minimum among all pairs of points in separate clusters ($|s_2 - s_1| = split_C(X)$). Let $w_l$ and $w_2$ be two points in $X$ such that the distance between them is maximum among all pairs of points in the same cluster.

For the points in $X$ to have a $\sigsep$ clustering in $\tX$, the following condition should hold in the worst case,

\[
 s_2 (1-\alpha) - s_1 (1+\alpha) > w_2 (1+\alpha) - w_1 (1 -\alpha)
\]
\[
\implies (s_2 - s_1) - \alpha (s_1 + s_2) > (w_2 - w_1) + \alpha (w_1 + w_2)
\]

\[
\implies \alpha < \frac{(s_2 - s_1) - (w_2 - w_1)}{4} = \frac{split_C(X) - width_C(X)}{4}
\]

This implies that if the original clusters are $\sigsep$ and $\alpha$ is small enough then the optimal cluster memberships over $X$ and $\tX$ will be identical. 








